% GENERATED FILE. Edit canonical lessons, then run npm run generate:books.
% canonical-source-hash: fnv1a64:40b65b375b3314f8
% canonical-lessons: PT-C308-susto, PT-C308-assustar, PT-C308-panico, PT-C308-desmaiar, PT-C308-torcer, PT-R308-revisao-301-308

\chapter{Que susto!}
\label{ch:pt-que-susto}
\hlchaptermodality{308}

\begin{chapteropening}
\textbf{By the end of this chapter:} \emph{I can talk about frights and panic, fainting, and twisting an ankle.}

Complete the last lesson of chapter 308: I can talk about frights and panic, fainting, and twisting an ankle.
\end{chapteropening}

\section[o susto]{o susto — the fright, the scare}

\label{lesson:PT-C308-susto}



\begin{quote}
Before the new one: say the Portuguese for the fracture, then the Portuguese for the plaster.
\end{quote}

\subsection*{You'll want to know: o susto}
\textbf{o susto} — \textquotedblleft{}the fright, the scare\textquotedblright{}.

\begin{quote}
Que susto! Pensei que estava sozinho em casa.
\end{quote}

\begin{quote}
\emph{What a fright! I thought I was alone in the house.}
\end{quote}

\emph{Pregar um susto a alguém} — to give someone a fright.

\begin{cousinweb}
The origin of \emph{susto} is not settled; EtymDB lists it only beside Spanish \emph{susto}.
\end{cousinweb}

\begin{grammarlens}[title={What you've built}]
One more word for telling, explaining and describing at B1.
\end{grammarlens}

\subsection*{Your turn}
\begin{itemize}
\raggedright
  \item \emph{Say it:} \emph{o susto}
  \item \emph{Say it:} \emph{o susto}, once more
  \item \emph{Say it:} a sentence of your own with \emph{o susto}
  \item \emph{Recall:} say the Portuguese for the fracture, then the Portuguese for the plaster, then say \emph{o susto} again
\end{itemize}

\begin{tcolorbox}[breakable,colback=teal!4,colframe=teal!35!black,title={Before you move on}]
What does o susto mean? (\textquotedblleft{}The fright, the scare\textquotedblright{}.) Say it once more.
\end{tcolorbox}
\section[assustar]{assustar — to frighten, to scare}

\label{lesson:PT-C308-assustar}



\begin{quote}
Before the new one: say the Portuguese for the plaster, then the Portuguese for the fright.
\end{quote}

\subsection*{You'll want to know: assustar}
\textbf{assustar} — \textquotedblleft{}to frighten, to scare\textquotedblright{}.

\begin{quote}
O barulho assustou as crianças.
\end{quote}

\begin{quote}
\emph{The noise frightened the children.}
\end{quote}

\emph{Assustar-se} — to get a fright.

\begin{cousinweb}
\emph{assustar} is built on \emph{susto} \textquotedblleft{}a fright\textquotedblright{}; EtymDB lists it beside Spanish \emph{asustar}, and the origin is not settled.
\end{cousinweb}

\begin{grammarlens}[title={What you've built}]
One more word for telling, explaining and describing at B1.
\end{grammarlens}

\subsection*{Your turn}
\begin{itemize}
\raggedright
  \item \emph{Say it:} \emph{assustar}
  \item \emph{Say it:} \emph{assustar}, once more
  \item \emph{Say it:} a sentence of your own with \emph{assustar}
  \item \emph{Recall:} say the Portuguese for the plaster, then the Portuguese for the fright, then say \emph{assustar} again
\end{itemize}

\begin{tcolorbox}[breakable,colback=teal!4,colframe=teal!35!black,title={Before you move on}]
What does assustar mean? (\textquotedblleft{}To frighten, to scare\textquotedblright{}.) Say it once more.
\end{tcolorbox}
\section[o pânico]{o pânico — panic}

\label{lesson:PT-C308-panico}



\begin{quote}
Before the new one: say the Portuguese for the fright, then the Portuguese for to frighten.
\end{quote}

\subsection*{You'll want to know: o pânico}
\textbf{o pânico} — \textquotedblleft{}panic\textquotedblright{}.

\begin{quote}
Houve pânico no centro comercial.
\end{quote}

\begin{quote}
\emph{There was panic in the shopping centre.}
\end{quote}

\emph{Entrar em pânico} — to panic.

\begin{cousinweb}
\emph{pânico} comes from Greek \emph{panikós} \textquotedblleft{}of Pan\textquotedblright{} (EtymDB), the god of the wild whose sudden cry was said to spread terror.
\end{cousinweb}

\begin{grammarlens}[title={What you've built}]
One more word for telling, explaining and describing at B1.
\end{grammarlens}

\subsection*{Your turn}
\begin{itemize}
\raggedright
  \item \emph{Say it:} \emph{o pânico}
  \item \emph{Say it:} \emph{o pânico}, once more
  \item \emph{Say it:} a sentence of your own with \emph{o pânico}
  \item \emph{Recall:} say the Portuguese for the fright, then the Portuguese for to frighten, then say \emph{o pânico} again
\end{itemize}

\begin{tcolorbox}[breakable,colback=teal!4,colframe=teal!35!black,title={Before you move on}]
What does o pânico mean? (\textquotedblleft{}Panic\textquotedblright{}.) Say it once more.
\end{tcolorbox}
\section[desmaiar]{desmaiar — to faint}

\label{lesson:PT-C308-desmaiar}



\begin{quote}
Before the new one: say the Portuguese for to frighten, then the Portuguese for panic.
\end{quote}

\subsection*{You'll want to know: desmaiar}
\textbf{desmaiar} — \textquotedblleft{}to faint\textquotedblright{}.

\begin{quote}
Com o calor, uma senhora desmaiou no autocarro.
\end{quote}

\begin{quote}
\emph{In the heat, a lady fainted on the bus.}
\end{quote}

\emph{O desmaio} — the faint.

\begin{cousinweb}
\emph{desmaiar} came from Old French \emph{desmaier} (EtymDB), from Vulgar Latin \emph{exmagāre} \textquotedblleft{}to rob of strength\textquotedblright{}.
\end{cousinweb}

\begin{grammarlens}[title={What you've built}]
One more word for telling, explaining and describing at B1.
\end{grammarlens}

\subsection*{Your turn}
\begin{itemize}
\raggedright
  \item \emph{Say it:} \emph{desmaiar}
  \item \emph{Say it:} \emph{desmaiar}, once more
  \item \emph{Say it:} a sentence of your own with \emph{desmaiar}
  \item \emph{Recall:} say the Portuguese for to frighten, then the Portuguese for panic, then say \emph{desmaiar} again
\end{itemize}

\begin{tcolorbox}[breakable,colback=teal!4,colframe=teal!35!black,title={Before you move on}]
What does desmaiar mean? (\textquotedblleft{}To faint\textquotedblright{}.) Say it once more.
\end{tcolorbox}
\section[torcer]{torcer — to twist, to sprain; to support (a team)}

\label{lesson:PT-C308-torcer}



\begin{quote}
Before the new one: say the Portuguese for panic, then the Portuguese for to faint.
\end{quote}

\subsection*{You'll want to know: torcer}
\textbf{torcer} — \textquotedblleft{}to twist, to sprain; to support (a team)\textquotedblright{}.

\begin{quote}
Torci o pé nas escadas.
\end{quote}

\begin{quote}
\emph{I twisted my ankle on the stairs.}
\end{quote}

\emph{Torcer por uma equipa} — to support a team.

\begin{cousinweb}
\emph{torcer} is inherited from Vulgar Latin \emph{torcēre}, from Latin \emph{torquēre} \textquotedblleft{}to twist\textquotedblright{} (EtymDB).
\end{cousinweb}

\begin{grammarlens}[title={What you've built}]
One more word for telling, explaining and describing at B1.
\end{grammarlens}

\subsection*{Your turn}
\begin{itemize}
\raggedright
  \item \emph{Say it:} \emph{torcer}
  \item \emph{Say it:} \emph{torcer}, once more
  \item \emph{Say it:} a sentence of your own with \emph{torcer}
  \item \emph{Recall:} say the Portuguese for panic, then the Portuguese for to faint, then say \emph{torcer} again
\end{itemize}

\begin{tcolorbox}[breakable,colback=teal!4,colframe=teal!35!black,title={Before you move on}]
What does torcer mean? (\textquotedblleft{}To twist, to sprain; to support (a team)\textquotedblright{}.) Say it once more.
\end{tcolorbox}
\section[(dialogue)]{Revisão dos capítulos 301–308}

\label{lesson:PT-R308-revisao-301-308}



\begin{quote}
Say the Portuguese for the fire and the explosion.
\end{quote}

\subsection*{Your turn}
Say these aloud:

\begin{itemize}
\raggedright
  \item the fire, the explosion, to explode, the fire extinguisher, the burn
  \item the quake, the earthquake, the tremor, to collapse, the rubble
  \item the flood, to flood, the drowning, the alert, to warn
  \item the catastrophe, the disaster, the tragedy, the landslide, to bury under earth or rubble
  \item the ambulance, the siren, to help, the first-aider, the stretcher
  \item to evacuate, to rescue, the rescue, to survive, the survivor
  \item the injured person, to injure, the injury, the fracture, the plaster
  \item the fright, to frighten, panic, to faint, to twist
\end{itemize}

\begin{tcolorbox}[breakable,colback=teal!4,colframe=teal!35!black,title={Before you move on}]
The fire? (\textbf{o incêndio}.) The explosion? (\textbf{a explosão}.) To explode? (\textbf{explodir}.) The fire extinguisher? (\textbf{o extintor}.) The burn? (\textbf{a queimadura}.) The quake? (\textbf{o sismo}.) The earthquake? (\textbf{o terramoto}.) The tremor? (\textbf{o tremor}.) To collapse? (\textbf{ruir}.) The rubble? (\textbf{os escombros}.) The flood? (\textbf{a inundação}.) To flood? (\textbf{inundar}.) The drowning? (\textbf{o afogamento}.) The alert? (\textbf{o alerta}.) To warn? (\textbf{alertar}.) The catastrophe? (\textbf{a catástrofe}.) The disaster? (\textbf{o desastre}.) The tragedy? (\textbf{a tragédia}.) The landslide? (\textbf{o deslizamento}.) To bury under earth or rubble? (\textbf{soterrar}.) The ambulance? (\textbf{a ambulância}.) The siren? (\textbf{a sirene}.) To help? (\textbf{socorrer}.) The first-aider? (\textbf{o socorrista}.) The stretcher? (\textbf{a maca}.) To evacuate? (\textbf{evacuar}.) To rescue? (\textbf{resgatar}.) The rescue? (\textbf{o resgate}.) To survive? (\textbf{sobreviver}.) The survivor? (\textbf{o sobrevivente}.) The injured person? (\textbf{o ferido}.) To injure? (\textbf{ferir}.) The injury? (\textbf{o ferimento}.) The fracture? (\textbf{a fratura}.) The plaster? (\textbf{o gesso}.) The fright? (\textbf{o susto}.) To frighten? (\textbf{assustar}.) Panic? (\textbf{o pânico}.) To faint? (\textbf{desmaiar}.) To twist? (\textbf{torcer}.)
\end{tcolorbox}
